\chapter{Verifica e validazione}
\label{chap:verifica-validazione}

Il presente capitolo descrive le attività di testing e valutazione effettuate al fine di verificare il corretto funzionamento delle funzionalità implementate 
durante l'attività di stage. Le modalità di verifica sono state definite in funzione della natura dei diversi componenti sviluppati, distinguendo tra funzionalità 
caratterizzate da un comportamento deterministico e componenti basate sull'utilizzo di modelli linguistici.
\newline
Per il sistema di generazione degli alert sono stati realizzati unit test volti a verificare individualmente le regole implementate, utilizzando dati controllati e simulati 
per valutare il corretto rilevamento delle condizioni previste. Per il modulo XAI è stata invece adottata una metodologia di valutazione automatizzata delle risposte generate 
dal modello linguistico, attraverso specifiche classi di \textit{eval} in grado di analizzare le risposte ottenute a partire da richieste note e verificarne il rispetto dei 
criteri definiti.
\newline
Le sezioni successive descrivono nel dettaglio le strategie adottate per le due componenti, illustrando le modalità di esecuzione dei test e i criteri utilizzati per valutarne 
i risultati.

\section{Test delle regole alert}
Il metodo utilizzato per verificare il corretto funzionamento delle regole alert è basato sull'utilizzo di \textit{Unit Test}. Nella cartella 
\texttt{tests/alerts} sono stati sviluppati undici test, uno per ciascuna regola implementata. Tutti i test seguono una struttura comune, in modo da garantire uniformità 
nella verifica delle diverse funzionalità.
\newline
La preparazione dei dati necessari all'esecuzione dei test avviene attraverso i seguenti passaggi:
\begin{itemize}
    \item utilizzo delle \textbf{factory}, definite nella cartella \texttt{database/factories}, che permettono di creare istanze dei modelli Eloquent utilizzando dati controllati e 
        specifici per l'ambiente di testing;
    \item creazione dei modelli necessari alla verifica della singola regola, includendo anche i parametri opzionali che possono essere configurati dall'utente;
    \item definizione di una matrice di test contenente i diversi casi da verificare, specificando per ciascuno i dati forniti alla regola e il risultato atteso dal metodo 
        \texttt{check}.
\end{itemize}

L'utilizzo della matrice di test permette di verificare, attraverso un'unica struttura, sia i casi in cui le condizioni della regola vengono soddisfatte e deve quindi essere 
rilevata una violazione, sia i casi in cui non deve essere generato alcun alert.
\newline
L'ambiente di testing è configurato tramite il file \texttt{phpunit.xml}, nel quale vengono definite le impostazioni utilizzate durante l'esecuzione dei test, tra cui il database 
di test, le modalità di connessione e le configurazioni relative agli altri servizi utilizzati dall'applicazione.

\begin{listing}[H]
\captionsetup{type=listing, belowskip=0pt}
\inputminted[fontsize=\small, breaklines, breakanywhere=true, linenos=true, bgcolor=gray!10]{php}{code/test.php}
\caption{Esempio di una funzione di test per la regola BuyOrdersRule}
\label{listing:php_test}
\end{listing}

\section{Test del modulo XAI}
Data la natura non deterministica dei modelli linguistici, il testing del modulo XAI presenta alcune difficoltà rispetto alla verifica delle funzionalità tradizionali. La stessa 
richiesta può infatti produrre risposte differenti tra un'esecuzione e l'altra, pur risultando corrette dal punto di vista del contenuto. Per questo motivo, non è possibile 
basare la valutazione esclusivamente sul confronto diretto tra la risposta ottenuta e una risposta attesa.
\newline
Per valutare il comportamento del modulo è stato quindi sviluppato il comando \texttt{php artisan ai:eval}, che ricerca i file JSON contenenti i casi di test e le relative 
richieste da sottoporre al modello. Ogni richiesta viene eseguita dieci volte, permettendo di valutare il comportamento del modulo su più generazioni della stessa richiesta e 
di tenere conto della variabilità delle risposte prodotte dall'LLM.
\newline
Per ciascun caso di test vengono definite delle aspettative relative al contenuto che il modello dovrebbe fornire. Poiché la stessa informazione può essere espressa utilizzando 
forme linguistiche differenti, è stata implementata la classe \texttt{ExpectationResolver}, che si occupa di normalizzare i valori utilizzati nella valutazione. In particolare, 
vengono gestiti valori quali date, numeri e stringhe, permettendo di confrontare il contenuto delle informazioni ottenute senza richiedere una corrispondenza lessicale esatta.
\newline
La classe \texttt{BehaviourGrader} valuta invece il comportamento del modello durante l'elaborazione della richiesta. In particolare, verifica i tool utilizzati, il loro ordine 
di invocazione, i modelli interessati dalle operazioni e le operazioni effettuate sul database. Questo permette di verificare non solo il risultato finale, ma anche che il 
modello abbia seguito il comportamento previsto per raggiungerlo.
\newline
La classe \texttt{AnswerGrader} si occupa della valutazione della risposta finale prodotta dal modello, verificandone il contenuto e la capacità di fornire all'utente una 
risposta coerente con le aspettative definite per il caso di test.
\newline
Infine, la classe \texttt{AiChatEvalCommand} coordina le diverse componenti di valutazione, eseguendo i casi di test e raccogliendo i risultati prodotti dai diversi 
\textit{grader}. Oltre alla correttezza della risposta e del comportamento del modello, vengono registrate ulteriori informazioni relative all'esecuzione, tra cui il tempo 
necessario per ottenere la risposta e il numero di token utilizzati. In questo modo è possibile ottenere una valutazione complessiva del comportamento del modulo XAI, 
considerando sia la correttezza delle risposte sia i costi computazionali associati alla loro generazione.


\section{Continuous Integration}
A completamento delle attività di testing, il progetto utilizza un processo di \textit{Continuous Integration (CI)} finalizzato a verificare automaticamente le modifiche 
introdotte nel codice. Ogni \textit{Pull Request} aperta sul repository viene sottoposta a una serie di controlli automatici prima di poter essere integrata nel ramo di 
\textit{develop}.
\newline
In particolare, la pipeline esegue \textit{PHPStan}, attraverso la configurazione adottata dal progetto, per effettuare l'analisi statica del codice e individuare possibili 
errori o comportamenti non conformi ai tipi e alle strutture previste. Viene inoltre eseguito \textit{PHP Insights}, utilizzato per verificare la qualità complessiva del codice 
e individuare eventuali problematiche relative alla sua struttura e manutenibilità.
\newline
La \textit{Pull Request} può essere integrata nel ramo di sviluppo solamente dopo il superamento di tutti i controlli previsti dalla pipeline. In questo modo, ogni modifica 
viene sottoposta automaticamente a una verifica prima di essere integrata, riducendo il rischio di introdurre codice non conforme agli standard adottati dal progetto e 
contribuendo a mantenere la qualità e la stabilità del codice.

\section{Validazione}
Per verificare la soddisfazione dei requisiti di qualità sono state svolte, periodicamente, attività di validazione in presenza del tutor aziendale, con l'obiettivo di valutare 
il comportamento dell'applicazione in scenari rappresentativi del reale dominio di utilizzo.
\newline
A tale scopo sono stati eseguiti diversi scenari utilizzando dati provenienti da dump forniti dall'azienda proponente, così da rendere le condizioni di prova il più possibile 
aderenti a quelle presenti nell'ambiente reale. Le attività di validazione hanno riguardato principalmente le prestazioni del motore di alert, i tempi di risposta e la coerenza 
delle risposte del modulo XAI e la correttezza delle informazioni registrate dal sistema di logging.

\subsection{Esecuzione del motore alert su dati reali}
È stato eseguito il comando \texttt{alert} dopo aver abilitato diverse regole alert per molteplici utenti, con l'obiettivo di valutare il comportamento del motore in presenza 
di una quantità significativa di dati. Durante l'esecuzione sono stati monitorati i tempi necessari per completare le verifiche delle regole, verificando che l'elaborazione non 
richiedesse tempi tali da compromettere il normale funzionamento del sistema.
\newline
La prova ha inoltre permesso di verificare la tempestività nella generazione e nell'invio delle notifiche, valutando quindi il rispetto dei requisiti relativi ai tempi di 
risposta del sistema di alert.

\subsection{Richiesta di spiegazioni al modulo XAI su dati reali}
Sono state effettuate molteplici richieste al modulo XAI utilizzando dati reali e simulando anche richieste concorrenti. L'obiettivo della prova è stato verificare che 
l'infrastruttura scelta per l'hosting del modello linguistico fosse in grado di gestire le richieste con tempi di risposta accettabili e che le spiegazioni prodotte risultassero 
coerenti con i dati effettivamente presenti nel sistema.
\newline
Le richieste sono state inoltre utilizzate per verificare il corretto funzionamento dell'interazione tra il modulo XAI, i tool disponibili e i dati del gestionale, assicurandosi 
che le informazioni riportate nelle risposte fossero riconducibili ai dati effettivamente recuperati dal sistema.

\subsection{Verifica del logging delle operazioni su dati reali}
Sono state eseguite molteplici operazioni di creazione, modifica ed eliminazione su dati già presenti nel sistema, utilizzando i dati forniti dagli stakeholder. Per 
ciascuna operazione è stata verificata la corretta registrazione delle informazioni all'interno del sistema di audit, controllando in particolare l'operazione effettuata, 
l'utente responsabile, la data e le modifiche apportate ai dati.
\newline
La prova ha inoltre permesso di verificare la corretta visualizzazione delle informazioni registrate e la facilità di consultazione dello storico delle operazioni, con 
l'obiettivo di garantire un'interfaccia semplice e intuitiva per l'analisi dei log.

\subsection{Risultati}
Le attività di validazione si sono rivelate fondamentali per individuare alcune criticità di utilizzo che non erano emerse durante le precedenti fasi di sviluppo e testing. 
A seguito dell'individuazione di tali problematiche sono state apportate le opportune correzioni e sono state ripetute le prove di validazione.
\newline
Al termine delle attività, tutti gli scenari di validazione hanno prodotto risultati conformi alle aspettative. In accordo con il tutor aziendale, è stata quindi confermata 
la corretta integrazione e l'adeguatezza dei moduli implementati rispetto al dominio di utilizzo dell'azienda committente.
\newpage

\section{Tracciamento dei requisiti}
Nelle successive tabelle vengono riportati gli stati di tutti i requisiti identificati durante l'attività di analisi dei requisiti.
\subsection{Requisti funzionali}
\keepXColumns
\rowcolors{2}{lightblue}{white}
    \begin{longtable}{|l|p{0.55\textwidth}|l|}
        \hline
        \rowcolor{gold}
        \textbf{Codice} & \textbf{Descrizione requisito} & \textbf{Soddisfatto} \\
        \hline
        \endhead

        \hline
        \rowcolor{white}
        \caption{Tabella riguardante lo stato dei requisiti funzionali del progetto alla fine del periodo di stage} \label{tab:requisiti_stato_funzionali} \\
        \endlastfoot 

        \textbf{RFN-1} & Configurazione del modulo XAI & Sì \\
        \hline
        \textbf{RFN-2} & Selezione del modello LLM per il modulo XAI & Sì \\
        \hline
        \textbf{RFN-3} & Configurazione degli utenti abilitati al modulo XAI & Sì \\
        \hline
        \textbf{RFN-4} & Configurazione dei tool disponibili per il modulo XAI & Sì \\
        \hline
        \textbf{RFN-5} & Richiesta di spiegazioni al modulo XAI su una decisione del sistema & Sì \\
        \hline
        \textbf{RFN-6} & Gestione della conversazione con il modulo XAI & Sì \\
        \hline
        \textbf{RFN-7} & Chiamata a tool per spiegazioni del modulo XAI & Sì \\
        \hline
        \textbf{RFN-8} & Modello LLM non raggiungibile & Sì \\
        \hline
        \textbf{RFN-9} & Limite token di risposta superato & Sì \\
        \hline
        \textbf{RFN-10} & Permessi non sufficienti per effettuare la richiesta & Sì \\
        \hline
        \textbf{RFN-11} & Richiesta creazione nuovi record & Sì \\
        \hline
        \textbf{RFN-12} & Chiamata al tool per creazione record del modulo XAI & Sì \\
        \hline
        \textbf{RFN-13} & Dati forniti incorretti per la creazione del record & Sì \\
        \hline
        \textbf{RFN-14} & Configurazione delle regole alert & Sì \\
        \hline
        \textbf{RFN-15} & Visualizzazione delle notifiche alert & Sì \\
        \hline
        \textbf{RFN-16} & Numero elevato di violazioni rilevate & Sì \\
        \hline
        \textbf{RFN-17} & Visualizzazione del dettaglio dell'alert & Sì \\
        \hline
        \textbf{RFN-18} & Visualizzazione del modello associato all'alert & Sì \\
        \hline
        \textbf{RFN-19} & Richiesta al modulo XAI di una spiegazione relativa all'alert generato & Sì \\
        \hline
        \textbf{RFN-20} & Visualizzazione log operazioni & Sì \\
        \hline
        \textbf{RFN-21} & Visualizzazione dei dettagli di un'operazione & Sì \\
        \hline
    \end{longtable}

\newpage

\subsection{Requisti di qualità}
\keepXColumns
\rowcolors{2}{lightblue}{white}
    \begin{longtable}{|l|p{0.55\textwidth}|l|}
        \hline
        \rowcolor{gold}
        \textbf{Codice} & \textbf{Descrizione requisito} & \textbf{Soddisfatto} \\
        \hline
        \endhead

        \hline
        \rowcolor{white}
        \caption{Tabella riguardante lo stato dei requisiti di qualità del progetto alla fine del periodo di stage} \label{tab:requisiti_stato_qualitativi} \\
        \endlastfoot 

        \textbf{RQN-1} & Il sistema deve garantire tempi di risposta adeguati nell'interazione con il modulo XAI, evitando attese non necessarie da parte dell'utente & Sì \\
        \hline
        \textbf{RQN-2} & Il sistema deve garantire la sicurezza delle informazioni, impedendo agli utenti di accedere a dati per i quali non possiedono i permessi necessari & Sì \\
        \hline
        \textbf{RQN-3} & Il sistema deve garantire la tracciabilità delle operazioni effettuate sui dati, registrando creazione, modifica ed eliminazione dei modelli & Sì \\
        \hline
        \textbf{RQN-4} & Le spiegazioni generate dal modulo XAI devono essere comprensibili per l'utente e presentate in linguaggio naturale & Sì \\
        \hline
        \textbf{RQN-5} & Il sistema deve garantire tempi di risposta adeguati durante il recupero dei dati e nella conseguente generazione e notifica degli alert & Sì \\
        \hline
        \textbf{RQN-6} & Il sistema deve notificare tempestivamente l'utente in caso di eccezioni nell'uso di modulo XAI e alert per ridurre le attese inutili & Sì \\
    \end{longtable}


\subsection{Requisti di vincolo}
\keepXColumns
\rowcolors{2}{lightblue}{white}
    \begin{longtable}{|l|p{0.55\textwidth}|l|}
        \hline
        \rowcolor{gold}
        \textbf{Codice} & \textbf{Descrizione requisito} & \textbf{Soddisfatto} \\
        \hline
        \endhead

        \hline
        \rowcolor{white}
        \caption{Tabella riguardante lo stato dei requisiti di vincolo del progetto alla fine del periodo di stage} \label{tab:requisiti_stato_vincolo} \\
        \endlastfoot 

        \textbf{RVN-1} & Il sistema deve essere sviluppato in PHP & Sì \\
        \hline
        \textbf{RVN-2} & L'interfaccia grafica deve essere sviluppata in Filament & Sì \\
        \hline
        \textbf{RVN-3} & Il modulo XAI deve fornire le spiegazioni generate in lingua italiana & Sì \\
        \hline
        \textbf{RVN-4} & L'integrazione con i modelli linguistici deve avvenire attraverso le API messe a disposizione dai relativi provider & Sì \\
        \hline
    \end{longtable}

Dall'attività di verifica finale è emerso che tutti i requisiti obbligatori individuati in fase di analisi sono stati soddisfatti con successo.




\newpage